% Provenance (never printed): see sources/outline_sources.md (Year 2: eccentricity, obliquity and
% precession; insolation computed from spherical astronomy; the Quaternary cycles; carbon dioxide
% and feedbacks). Orbital elements and tabulated insolation of Berger and Loutre (1991) via NOAA
% NCEI (data/derived/bachelor-2-orbit91, fetch_data Book 4 block); periods after Laskar et al.
% (2004) and the NASA Milankovitch page; history after the NASA Earth Observatory feature (2000);
% terminations after Barker and Knorr (2021); the mid-Pleistocene transition after Clark et al.
% (2006) and Koehler and van de Wal (2020). Figures from figdata/bachelor-2/orbital.py; values of
% the exercises and problem computed in the Book 4 scratch calculator calc/ch25.py.
\chapter{Orbitale forcering en de ijstijden}\label{ch:b2:orbital-forcing}

Over het laatste miljoen jaar heeft de vorm van de baan van de aarde het zonlicht dat de hele
planeet in een jaar ontvangt met minder dan twee duizendsten veranderd: tussen \num{340.00} en
\qty{340.54}{W.m^{-2}} aan de top van de atmosfeer. Over hetzelfde miljoen jaar schommelde de daggemiddelde
instraling van half juli op \ang{65}~NB tussen ongeveer 388 en \qty{500}{W.m^{-2}}, met meer dan
een kwart. De ijstijden volgden het tweede getal, niet het eerste. Dit hoofdstuk berekent beide
uit de elementen van de baan van de aarde, en vraagt hoe het zonlicht op één breedte in één seizoen
het tempo kan aangeven van de groei en het verval van ijskappen die continenten bedekken.
% ledger: data:b4-orbital, berger91:readme

\begin{recall}
\textbf{Uit de natuurkundeboeken}: de ellips, haar halve lange as $a$ en excentriciteit $e$, de wetten
van Kepler. \textbf{Uit Boek~3}: de ecliptica en de scheefstand $\varepsilon$, de
declinatie van de zon $\delta$ met $\sin\delta = \sin\varepsilon\sin\lambda$ ($\lambda$ de ware
lengte van de zon, langs de ecliptica geteld vanaf het lentepunt), het tropisch jaar, de daggemiddelde
instraling
\[
  \bar Q = \frac{S_0}{\pi}\left(\frac{a}{r}\right)^{2}\bigl(h_0\sin\varphi\sin\delta +
  \cos\varphi\cos\delta\sin h_0\bigr),\qquad \cos h_0 = -\tan\varphi\tan\delta ,
\]
en haar globaal en jaarlijks gemiddelde, $S_0/(4\sqrt{1 - e^2})$. \textbf{Uit de schoolboeken}: de
glaciale cyclus en het Laatste Glaciale Maximum; stralingsforcering. \textbf{Uit eerdere hoofdstukken}: de
\omterm{def:b2:precambrian-earth:ice-albedo}{ijs--albedoterugkoppeling} (\cref{def:b2:precambrian-earth:ice-albedo}); de benthische stapelreeks, de reeks van EPICA
Dome C en ouderdomsmodellen (\cref{ch:b2:climate-proxies}).
\end{recall}
% ledger: earth:TSI, earth:obliquity, earth:ecc

\section{De baan en haar drie elementen}\label{sec:b2:orbital-forcing:elements}

De baan van de aarde is een ellips met de zon in één brandpunt. De planeten, vooral Jupiter en Saturnus,
vervormen haar langzaam en draaien haar in haar vlak; de aantrekking van de zon en de maan op de equatoriale
uitstulping doet de rotatieas tollen als een tol. Drie grootheden beschrijven wat er voor het
zonlicht toe doet: hoe langgerekt de ellips is, hoe scheef de as staat, en waar langs de baan de
seizoenen vallen.
% ledger: nasa:milank-causes

\begin{definition}[Lengte van het perihelium; klimatologische precessie]\label{def:b2:orbital-forcing:elements}
De \emph{lengte van het perihelium}\index{lengte van het perihelium} $\varpi$ is de hoek, gemeten in
het baanvlak vanaf de richting van het lentepunt, tussen die richting en de
richting van het perihelium gezien vanaf de zon; ze wordt geteld vanaf het bewegende lentepunt, dus ze
verandert zowel door het tollen van de as als door het draaien van de baan. De \emph{klimatologische
precessie}\index{klimatologische precessie} is de grootheid $e\sin\varpi$, ook
precessieindex genoemd.
\end{definition}
% ledger: berger91:readme

\begin{omfigure}
\centering
\begin{tikzpicture}[scale=0.82, font=\scriptsize, >=Stealth]
\draw[thick, rotate=101.37] (-0.72,0) ellipse (2.4 and 2.2895);
\draw[densely dashed, black!60] (-0.331,1.647) -- (0.615,-3.059);
\draw[->, black!70] (0,0) -- (3.9,0);
\node[anchor=south, align=center] at (3.75,0.06) {$\gamma$, naar het\\lentepunt};
\draw[->, omDef] (0.75,0) arc[start angle=0, end angle=101.37, radius=0.75];
\node[omDef] at (0.78,0.62) {$\varpi$};
\fill[ocRed] (0,0) circle (0.13);
\node[anchor=north west] at (0.05,-0.08) {zon};
\foreach \x/\y in {-2.062/0, 0/-3.094, 2.321/0, 0/1.688} {\fill[ocBlue] (\x,\y) circle (0.08);}
\node[anchor=east, align=right] at (-2.15,0) {maart-\\equinox};
\node[anchor=north west, align=left] at (2.38,-0.08) {september-\\equinox};
\node[anchor=west] at (0.1,1.92) {decemberzonnewende};
\node[anchor=east] at (-0.12,-3.25) {junizonnewende};
\fill[black] (-0.331,1.647) circle (0.06);
\node[anchor=east] at (-0.42,1.72) {perihelium};
\fill[black] (0.615,-3.059) circle (0.06);
\node[anchor=west] at (0.7,-3.15) {aphelium};
\draw[->, thick, black!60] (-1.9,-1.75) arc[start angle=220, end angle=245, radius=2.6];
\end{tikzpicture}
\par\smallskip{\small De baan van de aarde gezien vanaf de noordkant, met de excentriciteit overdreven
tot 0.3 en de huidige \omterm{def:b2:orbital-forcing:elements}{lengte van het perihelium}, $\varpi = \ang{101.4}$. De aarde beweegt tegen de klok in
(pijl); de punten geven haar positie bij de equinoxen en zonnewenden aan. Het perihelium valt enkele
dagen na de decemberzonnewende, het aphelium enkele dagen na de junizonnewende. Schematisch.}
\end{omfigure}
% ledger: berger91:rows, nasa:milank-detail

Gezien vanaf de aarde beweegt de zon langs de ecliptica, en haar ware lengte $\lambda$ is de
hoek vanaf het lentepunt: $\lambda = \ang{90}$ bij de junizonnewende, \ang{270} bij de
decemberzonnewende. De aarde staat in het perihelium wanneer de zon in de richting tegenover het perihelium staat,
bij $\lambda = \varpi + \ang{180}$; vandaag is $\varpi = \ang{101.37}$ en valt het perihelium bij
$\lambda = \ang{281.37}$, ongeveer elf dagen na de decemberzonnewende, rond 3 januari.
% ledger: berger91:rows, nasa:milank-detail

\begin{definition}[Milankovitch-parameters]\label{def:b2:orbital-forcing:milankovitch}
De \emph{Milankovitch-parameters}\index{Milankovitch-parameters} zijn de drie
baangrootheden die de verdeling van het zonlicht over de aarde door de seizoenen en de
breedten bepalen: de excentriciteit $e$, de scheefstand $\varepsilon$ en de \omterm{def:b2:orbital-forcing:elements}{klimatologische precessie}
$e\sin\varpi$. De \emph{Milankovitch-theorie}\index{Milankovitch-theorie} stelt dat hun trage
variaties het tempo aangeven van de groei en het verval van de kwartaire ijskappen, via de zomerinstraling
op hoge noordelijke breedten.
\end{definition}
% ledger: nasa:milank-detail, eo00:milankovitch

Elke parameter varieert met zijn eigen ritmen. De excentriciteit heeft hoofdperioden bij \num{95000},
\num{99000}, \num{124000} en \num{131000} jaar, die samen een cyclus van ongeveer
\num{100000} jaar vormen, en een periode van \num{405000} jaar, de stabielste van alle; over het laatste
miljoen jaar varieerde ze van \num{0.0046} tot \num{0.0565}, en vandaag, met \num{0.0172}, neemt ze langzaam
af. De scheefstand varieert met een hoofdperiode van \num{40978} jaar, tussen \ang{22.1} en
\ang{24.5} over het laatste miljoen jaar; met \ang{23.44} neemt ze af, voorbij haar laatste maximum van
ongeveer \num{10000} jaar geleden. De as tolt rond in ongeveer \num{25770} jaar, terwijl de
baan in ongeveer \num{112000} jaar draait; samen verplaatsen ze het perihelium door de seizoenen met
hoofdperioden bij \num{19000}, \num{22400} en \num{23700} jaar.
% ledger: la04:eccentricity, la04:obliquity, la04:precession, nasa:milank-detail, data:b4-orbital

\begin{omfigure}
\centering
% figdata/bachelor-2/orbital.py part elements
\begin{tikzpicture}
\begin{groupplot}[group style={group size=1 by 3, vertical sep=0.3cm, xlabels at=edge bottom,
  xticklabels at=edge bottom}, width=0.96\linewidth, height=3.1cm, xmin=0, xmax=1000,
  xtick={0,100,...,1000}, tick label style={font=\scriptsize}, label style={font=\small},
  scaled y ticks=false, yticklabel style={/pgf/number format/fixed,
  /pgf/number format/precision=2}]
\nextgroupplot[ylabel={$e$}, ymin=0, ymax=0.085, ytick={0,0.02,0.04,0.06}]
\addplot[omDef, thin] table[x=age_ka, y=ecc] {figdata/out/bachelor-2-orbital-elements.dat};
\node[font=\scriptsize, anchor=north west] at (rel axis cs:0.01,0.97) {a\enspace excentriciteit};
\nextgroupplot[ylabel={$\varepsilon$ (\unit{\degree})}, ymin=21.9, ymax=25.4, ytick={22,23,24}]
\addplot[ocBlue, thin] table[x=age_ka, y=obl] {figdata/out/bachelor-2-orbital-elements.dat};
\node[font=\scriptsize, anchor=north west] at (rel axis cs:0.01,0.97) {b\enspace scheefstand};
\nextgroupplot[ylabel={$e\sin\varpi$}, ymin=-0.07, ymax=0.12, ytick={-0.05,0,0.05},
  xlabel={ouderdom (duizenden jaren vóór 1950)}]
\addplot[ocRed, thin] table[x=age_ka, y=prec] {figdata/out/bachelor-2-orbital-elements.dat};
\node[font=\scriptsize, anchor=north west] at (rel axis cs:0.01,0.97) {c\enspace klimatologische
precessie};
\end{groupplot}
\end{tikzpicture}
\par\smallskip{\small De drie \omterm{def:b2:orbital-forcing:milankovitch}{Milankovitch-parameters} over het laatste miljoen jaar: (a) de
excentriciteit; (b) de scheefstand; (c) de \omterm{def:b2:orbital-forcing:elements}{klimatologische precessie}, waarvan de omhullende de excentriciteit is.
Gegevens: Berger en Loutre (1991), NOAA NCEI Paleoclimatology, opgehaald op 23 oktober 2026.}
\end{omfigure}
% ledger: data:b4-orbit91

De \omterm{def:b2:orbital-forcing:elements}{klimatologische precessie} schommelt snel, maar slechts zo ver als de excentriciteit toelaat: wanneer de baan
bijna cirkelvormig is, zoals \num{400000} jaar geleden, maakt het nauwelijks uit waar het perihelium valt. De
figuur toont het duidelijk: de omhullende van paneel (c) is paneel (a).
% ledger: data:b4-orbit91

\section{Zonlicht uit de baan}\label{sec:b2:orbital-forcing:insolation}

De daggemiddelde instraling die hierboven werd herhaald, heeft nog één ingrediënt nodig: de afstand tot de zon op een
gegeven dag.

\begin{proposition}[Afstand en de lengte van de zon]\label{prop:b2:orbital-forcing:distance}
Wanneer de ware lengte van de zon $\lambda$ is, is de afstand aarde--zon
\[
  \frac{r}{a} = \frac{1 - e^2}{1 - e\cos(\lambda - \varpi)} .
\]
Bij de junizonnewende ($\lambda = \ang{90}$) geldt, in eerste orde in $e$,
$(a/r)^2 \approx 1 - 2e\sin\varpi$; bij de decemberzonnewende $(a/r)^2 \approx 1 + 2e\sin\varpi$.
\end{proposition}

\begin{proof}
Op een ellips met de zon in een brandpunt is $r/a = (1 - e^2)/(1 + e\cos v)$, waarbij de ware anomalie $v$
de hoek bij de zon is van het perihelium naar de aarde. De aarde wordt vanaf de zon gezien in de
richting tegenover die van de zon gezien vanaf de aarde, dus haar lengte gezien vanaf de zon is
$\lambda + \ang{180}$, en $v = \lambda + \ang{180} - \varpi$, $\cos v = -\cos(\lambda - \varpi)$.
Bij $\lambda = \ang{90}$ is $\cos(\lambda - \varpi) = \sin\varpi$ en $(a/r)^2 = (1 -
e\sin\varpi)^2/(1 - e^2)^2 = 1 - 2e\sin\varpi + O(e^2)$; bij $\lambda = \ang{270}$ keert het teken om.
\end{proof}

De \omterm{def:b2:orbital-forcing:elements}{klimatologische precessie} is daarom de juiste index: ze meet hoeveel dichter bij de zon, of
verder weg, de aarde bij de junizonnewende staat. Wanneer $e\sin\varpi$ groot en positief is, zoals vandaag en
veel sterker \num{116000} jaar geleden, valt de noordelijke zomer bij het aphelium en is ze zwak; wanneer ze
groot en negatief is, zoals ongeveer \num{10000} jaar geleden, valt de noordelijke zomer bij het perihelium.
% ledger: berger91:rows

\begin{example}[De junizon op \ang{65}~NB, vandaag]\label{ex:b2:orbital-forcing:today}
Vandaag is $e = 0.017236$, $\varpi = \ang{101.37}$, $\varepsilon = \ang{23.446}$, en de tabellen gebruiken
$S_0 = \qty{1360}{W.m^{-2}}$. Bij de junizonnewende is $\delta = \varepsilon$; bij $\varphi =
\ang{65}$ is $\cos h_0 = -\tan\ang{65}\tan\ang{23.446} = -0.9301$, dus $h_0 = \ang{158.44} =
\qty{2.7654}{rad}$: de zon staat $2 \times 158.44/360 \times 24 = 21.1$~uur op. De haakjesterm van
de instralingsformule is $2.7654 \times 0.9063 \times 0.3979 + 0.4226 \times 0.9174 \times 0.3674 =
1.1397$, en op de gemiddelde afstand is $\bar Q = 432.90 \times 1.1397 = \qty{493.4}{W.m^{-2}}$. De
afstandsfactor is $(a/r)^2 = 0.9671$, dus $\bar Q = \qty{477.1}{W.m^{-2}}$.
\end{example}
% ledger: berger91:rows, berger91:readme, data:b4-orbital

De tweede wet van Kepler voegt een wending toe. Bij het aphelium beweegt de aarde langzamer: het seizoen dat
daar valt, is zwakker maar duurt langer. Noordelijke zomers zijn vandaag ongeveer 4.5 dagen langer dan noordelijke
winters, en lentes ongeveer drie dagen langer dan herfsten.
% ledger: nasa:milank-detail

\begin{proposition}[Jaargemiddelden hangen niet van de precessie af]\label{prop:b2:orbital-forcing:annual}
De energie die per oppervlakte-eenheid aan de top van de atmosfeer wordt ontvangen tussen twee gegeven lengten van
de zon, hangt niet af van $\varpi$, en de jaargemiddelde instraling op elke breedte is die van een
cirkelvormige baan met straal $a$, vermenigvuldigd met $1/\sqrt{1 - e^2}$. Op een pool is het jaargemiddelde
\[
  \bar Q_{\text{pool}} = \frac{S_0\sin\varepsilon}{\pi\sqrt{1 - e^2}} ,
\]
dat met de scheefstand toeneemt.
\end{proposition}

\begin{proof}
Volgens de tweede wet van Kepler bestrijkt de voerstraal oppervlakte met de constante snelheid $\pi ab/T$, en de oppervlakte die wordt bestreken
terwijl de aarde over $\mathrm d\lambda$ draait, is $\frac12 r^2\,\mathrm d\lambda$, dus $\mathrm dt =
T r^2\,\mathrm d\lambda/(2\pi ab)$, met $b = a\sqrt{1 - e^2}$. De instraling op een gegeven breedte
en lengte van de zon is $(a/r)^2$ keer haar waarde $Q_a(\lambda)$ op de afstand $a$, dus de
ontvangen energie is $\int Q_a(\lambda)(a/r)^2\,\mathrm dt = \frac{T}{2\pi\sqrt{1 - e^2}}\int
Q_a(\lambda)\,\mathrm d\lambda$: de afstand is weggevallen, en daarmee $\varpi$. Over het hele
jaar doorloopt $\lambda$ \ang{360} en is het gemiddelde $\frac{1}{2\pi\sqrt{1-e^2}}\int_0^{2\pi}
Q_a\,\mathrm d\lambda$. Op de noordpool is $Q_a = S_0\sin\delta = S_0\sin\varepsilon\sin\lambda$
voor $0 < \lambda < \pi$ en anders 0; de integraal is $2S_0\sin\varepsilon$.
\end{proof}
% ledger: earth:TSI

Precessie verplaatst zonlicht van het ene seizoen naar het andere, nooit een jaar in of uit; excentriciteit
verandert het jaartotaal met een factor $1/\sqrt{1 - e^2}$, hoogstens \qty{0.16}{\percent} over het
laatste miljoen jaar; scheefstand verplaatst zonlicht tussen breedten, verhoogt het jaargemiddelde aan de
polen en verlaagt het licht aan de evenaar. Alleen het patroon over de seizoenen en de breedten verandert
sterk, en dat toont de figuur hieronder.
% ledger: data:b4-orbital, nasa:milank-detail

\begin{omfigure}
\centering
% figdata/bachelor-2/orbital.py part lat
\begin{tikzpicture}
\begin{groupplot}[group style={group size=2 by 1, horizontal sep=1.6cm},
  tick label style={font=\scriptsize}, label style={font=\small}, height=5.4cm,
  width=0.48\linewidth, xmin=-90, xmax=90, xtick={-90,-60,...,90}, xlabel={breedte
  (\unit{\degree})}, ylabel={$\bar Q$ (\unit{W.m^{-2}})}]
\nextgroupplot[ymin=0, ymax=650, title={a\enspace junizonnewende}, title style={font=\scriptsize},
  legend style={font=\scriptsize, at={(0.03,0.97)}, anchor=north west, draw=none},
  legend cell align=left]
\addplot[black, thick] table[x=lat, y=jun_0] {figdata/out/bachelor-2-orbital-lat.dat};
\addlegendentry{vandaag}
\addplot[ocRed, thick, densely dashed] table[x=lat, y=jun_10] {figdata/out/bachelor-2-orbital-lat.dat};
\addlegendentry{\num{10000} jaar geleden}
\nextgroupplot[ymin=100, ymax=500, title={b\enspace jaargemiddelde}, title style={font=\scriptsize},
  legend style={font=\scriptsize, at={(0.5,0.03)}, anchor=south, draw=none},
  legend cell align=left]
\addplot[ocBlue, thick] table[x=lat, y=ann_lo] {figdata/out/bachelor-2-orbital-lat.dat};
\addlegendentry{$\varepsilon = \ang{22.1}$}
\addplot[ocRed, thick, densely dashed] table[x=lat, y=ann_hi] {figdata/out/bachelor-2-orbital-lat.dat};
\addlegendentry{$\varepsilon = \ang{24.5}$}
\end{groupplot}
\end{tikzpicture}
\par\smallskip{\small (a) Daggemiddelde instraling bij de junizonnewende, vandaag en \num{10000} jaar
geleden, toen de noordelijke zomer bij het perihelium viel en de scheefstand groter was: ongeveer
\qty{50}{W.m^{-2}} meer op de hoge noordelijke breedten. (b) Jaargemiddelde instraling voor de twee
uiterste scheefstanden van het laatste miljoen jaar, met de huidige excentriciteit: een grotere helling stuurt
meer zonlicht naar de polen en iets minder naar de tropen. Berekend uit de baanelementen
met $S_0 = \qty{1360}{W.m^{-2}}$.}
\end{omfigure}
% ledger: data:b4-orbital, data:b4-orbit91

\section{De instralingskromme berekenen}\label{sec:b2:orbital-forcing:curve}

De instralingskromme die Milankovitch opstelde, is nu een routineberekening. De baanelementen van
het verleden worden berekend door de bewegingen van de planeten te integreren; om de duizend jaar getabelleerd,
vormen ze de invoer. Voor elke ouderdom geeft de methode hieronder de instraling op elke breedte en in elk seizoen.

\begin{method}[Instraling van het verleden]\label{met:b2:orbital-forcing:curve}
\begin{enumerate}
  \item Lees $e$, $\varpi$ en $\varepsilon$ af voor de gewenste ouderdom.
  \item Kies het seizoen met de ware lengte van de zon $\lambda$: \ang{90} voor de junizonnewende;
        half juli in de tabellen komt overeen met $\lambda = \ang{120}$.
  \item Bereken $\delta$ uit $\sin\delta = \sin\varepsilon\sin\lambda$, en de uurhoek van de
        zonsondergang $h_0$ op de breedte $\varphi$ (pooldag of poolnacht als $|\tan\varphi\tan\delta| > 1$).
  \item Bereken de afstandsfactor $(a/r)^2$ met \cref{prop:b2:orbital-forcing:distance}.
  \item Vermenigvuldig: $\bar Q = (S_0/\pi)(a/r)^2(h_0\sin\varphi\sin\delta +
        \cos\varphi\cos\delta\sin h_0)$.
\end{enumerate}
\end{method}
% ledger: berger91:readme, data:b4-orbital

Toegepast op de 1001 ouderdommen van het laatste miljoen jaar reproduceert de methode de getabelleerde instraling van half juli
op \ang{65}~NB tot op \qty{0.012}{W.m^{-2}}, de afronding van de tabellen, en de
getabelleerde waarden voor juni van vandaag exact. De tabellen leggen de oorsprong van de tijd in 1950.
% ledger: data:b4-orbital, berger91:readme

\begin{omfigure}
\centering
% figdata/bachelor-2/orbital.py part jul65
\begin{tikzpicture}
\begin{groupplot}[group style={group size=1 by 2, vertical sep=0.3cm, xlabels at=edge bottom,
  xticklabels at=edge bottom}, width=0.96\linewidth, xmin=0, xmax=400, xtick={0,50,...,400},
  tick label style={font=\scriptsize}, label style={font=\small}]
\nextgroupplot[height=5.0cm, ymin=370, ymax=590, ylabel={$\bar Q$ op \ang{65}~NB
  (\unit{W.m^{-2}})}, legend style={font=\scriptsize, at={(0.5,1.02)}, anchor=south, draw=none,
  legend columns=3, /tikz/every even column/.append style={column sep=6pt}},
  legend cell align=left]
\addplot[black!45, line width=2.2pt] table[x=age_ka, y=tab] {figdata/out/bachelor-2-orbital-jul65.dat};
\addlegendentry{half juli, getabelleerd}
\addplot[ocRed, thin] table[x=age_ka, y=calc] {figdata/out/bachelor-2-orbital-jul65.dat};
\addlegendentry{half juli, berekend}
\addplot[ocBlue, thin, densely dashed] table[x=age_ka, y=jun65] {figdata/out/bachelor-2-orbital-jul65.dat};
\addlegendentry{junizonnewende, berekend}
\nextgroupplot[height=2.8cm, ymin=-0.02, ymax=0.02, ytick={-0.01,0,0.01},
  scaled y ticks=false, yticklabel style={/pgf/number format/fixed}, ylabel={residu},
  xlabel={ouderdom (duizenden jaren vóór 1950)}]
\addplot[only marks, mark=*, mark size=0.35pt, black] table[x=age_ka, y=resid]
  {figdata/out/bachelor-2-orbital-jul65.dat};
\end{groupplot}
\end{tikzpicture}
\par\smallskip{\small Boven: daggemiddelde instraling op \ang{65}~NB over de laatste \num{400000} jaar,
half juli zoals getabelleerd door Berger en Loutre (grijs, breed) en zoals berekend met de methode hierboven (rood),
en bij de junizonnewende (blauw, gestreept). Onder: berekende min getabelleerde waarden van half juli, in
\unit{W.m^{-2}}.}
\end{omfigure}
% ledger: data:b4-orbit91, data:b4-orbital

Twee kenmerken vallen op. De kromme wordt beheerst door de \omterm{def:b2:orbital-forcing:elements}{klimatologische precessie}, met pieken die zo'n
\num{20000} jaar uit elkaar liggen en waarvan de hoogte aanzwelt en afneemt met de excentriciteit; de scheefstand voegt
een tragere modulatie toe. En de schommelingen zijn groot: over de laatste \num{800000} jaar varieerde de waarde van half juli
van 388 tot \qty{497}{W.m^{-2}}.
% ledger: data:b4-orbital

\section{De Milankovitch-theorie en de glaciale cycli}\label{sec:b2:orbital-forcing:theory}

\begin{history}[Zonlicht op vijfenzestig graden noord]
Milutin Milankovitch, geboren in 1879, bouwde een wiskundige klimaattheorie op uit de variaties van het zonlicht met seizoen en breedte. Door de drie baanvariaties te combineren,
berekende hij de instraling van de breedten en de bijbehorende oppervlaktetemperaturen voor
\num{600000} jaar vóór 1800. Op voorstel van de klimatoloog Köppen koos hij de
zomerinstraling op \ang{65}~NB als de beslissende grootheid: de grote ijskappen waren dicht bij die
breedte gegroeid, en koele zomers smelten minder van de wintersneeuw, zodat sneeuw zich van jaar tot
jaar ophoopt. Hij achtte de scheefstand de belangrijkste van de drie parameters. Zo'n vijftig jaar lang
werd de theorie grotendeels genegeerd; Milankovitch stierf in 1958. In 1976 vonden Hays, Imbrie en Shackleton
in diepzeesedimenten van de laatste \num{450000} jaar spectrale pieken van klimaatvariabiliteit
bij ongeveer \num{23000}, \num{42000} en \num{100000} jaar, de perioden van precessie, scheefstand en
excentriciteit, en ze concludeerden dat veranderingen in de baan van de aarde de fundamentele oorzaak zijn van de
kwartaire ijstijden.
\end{history}
% ledger: eo00:milankovitch, nasa:milank-detail, hays76:periods

De figuur hieronder herhaalt hun test op de reeksen van \cref{ch:b2:climate-proxies}. Ze vergelijkt de
amplitudespectra, de sterkte van elke periode, van de juli-instraling op \ang{65}~NB en van de
benthische stapelreeks, in twee vensters van \num{700000} jaar.

\begin{omfigure}
\centering
% figdata/bachelor-2/orbital.py part spectrum
\begin{tikzpicture}
\begin{axis}[width=0.9\linewidth, height=5.6cm, xmin=10, xmax=150, ymin=0, ymax=1.08,
  xlabel={periode (duizenden jaren)}, ylabel={relatieve amplitude},
  xtick={19,23,41,100,150}, tick label style={font=\scriptsize}, label style={font=\small},
  legend style={font=\scriptsize, at={(0.5,1.02)}, anchor=south, draw=none, legend columns=3,
  /tikz/every even column/.append style={column sep=6pt}}, legend cell align=left]
\addplot[black!50, thick] table[x=period, y=amp_ins] {figdata/out/bachelor-2-orbital-spectrum.dat};
\addlegendentry{instraling \ang{65}~NB, 0--\qty{700}{ka}}
\addplot[ocRed, thick] table[x=period, y=amp_young] {figdata/out/bachelor-2-orbital-spectrum.dat};
\addlegendentry{LR04, 0--\qty{700}{ka}}
\addplot[ocBlue, thick, densely dashed] table[x=period, y=amp_old] {figdata/out/bachelor-2-orbital-spectrum.dat};
\addlegendentry{LR04, van \num{1250} tot \num{1950}~ka}
\end{axis}
\end{tikzpicture}
\par\smallskip{\small Amplitudespectra, elk geschaald op zijn eigen maximum, van de instraling van half juli
op \ang{65}~NB over de laatste \num{700000} jaar en van de benthische stapelreeks LR04 over hetzelfde venster
en over \num{1250000} tot \num{1950000} jaar geleden (trend verwijderd, getaperd). De instraling heeft haar
vermogen bij de precessieperioden, rond 19 tot 24 duizend jaar, en bij de periode van de scheefstand; het
recente ijsvolume piekt rond \num{100000} jaar, waar de instraling bijna niets heeft; het oudere
ijsvolume piekt bij \num{41000} jaar.}
\end{omfigure}
% ledger: data:b4-orbital, data:lr04, hays76:periods

\subsection*{Glacialen, interglacialen en terminaties}

\begin{definition}[Glaciale periode; interglaciaal; glaciale terminatie]\label{def:b2:orbital-forcing:glacials}
In de kwartaire reeks is een \emph{glaciale periode}\index{glaciale periode} een interval waarin grote
ijskappen de noordelijke continenten bedekken, vastgelegd door een hoge benthische $\delta^{18}$O en een laag zeeniveau;
een \emph{interglaciaal}\index{interglaciaal} is een interval met weinig landijs op het noordelijk
halfrond buiten Groenland, zoals nu. Een \emph{glaciale
terminatie}\index{glaciale terminatie} is een van de grote overgangen die een glaciale cyclus beëindigen,
van volledig glaciale naar interglaciale omstandigheden, veel sneller dan de groei die eraan voorafging.
\end{definition}
% ledger: kohler20:mpt, barker21:terminations, ipcc:ref-periods

Over het laatste half miljoen jaar ligt de dominante periode van het klimaat dicht bij \num{100000} jaar,
hoewel de forcering met die periode, die alleen van de excentriciteit komt, verwaarloosbaar is
naast die van scheefstand en precessie. De cycli zijn ook scheef: de benthische $\delta^{18}$O
registreert een stapsgewijze opbouw van ijs over ongeveer \num{90000} jaar, daarna een afbraak in ongeveer
\num{10000}, een zaagtand waarvan de dalende flanken de terminaties zijn. Het laatste \omterm{def:b2:orbital-forcing:glacials}{interglaciaal}, van
ongeveer \num{129000} tot \num{116000} jaar geleden, eindigde toen de zomerinstraling op \ang{65}~NB daalde
tot een van haar laagste waarden van de laatste \num{800000} jaar: de waarde bij de junizonnewende bereikte haar
minimum, \qty{438.6}{W.m^{-2}}, \num{116000} jaar geleden, en de waarde van half juli haar op een na laagste
minimum, \qty{392.2}{W.m^{-2}}, \num{114000} jaar geleden. Het Laatste Glaciale Maximum kwam zo'n
\num{95000} jaar later.
% ledger: barker21:terminations, ipcc:ref-periods, data:b4-orbital, data:b4-orbital-extremes

Hoe snel kan een ijskap haar forcering volgen? Het eenvoudigste model behandelt het ijsvolume $V$ als
relaxerend naar het volume $V_{\mathrm{eq}}$ dat de huidige instraling in stand zou houden, met een
responstijd $\tau$:
\[
  \frac{\mathrm dV}{\mathrm dt} = \frac{V_{\mathrm{eq}}(t) - V}{\tau} .
\]

\begin{proposition}[Een vertraagde respons]\label{prop:b2:orbital-forcing:lag}
Voor een forcering $V_{\mathrm{eq}} = A\sin\omega t$ nadert de oplossing van het vertragingsmodel
\[
  V = \frac{A}{\sqrt{1 + \omega^2\tau^2}}\,\sin\bigl(\omega(t - t_{\ell})\bigr),\qquad
  t_{\ell} = \frac{\arctan(\omega\tau)}{\omega} :
\]
het ijs volgt trage forcering bijna volledig en snelle forcering maar zwak, altijd te laat. De respons
op een som van sinusoïden is de som van de responsen: het model schept geen periode die zijn forcering
mist.
\end{proposition}

\begin{proof}
Probeer $V = B\sin\omega t + C\cos\omega t$: invullen geeft $\omega B\tau = -C$ en $-\omega C\tau = A -
B$, dus $B = A/(1 + \omega^2\tau^2)$, $C = -\omega\tau A/(1 + \omega^2\tau^2)$, een amplitude
$\sqrt{B^2 + C^2} = A/\sqrt{1 + \omega^2\tau^2}$ en een faseachterstand $\arctan(\omega\tau)$. Elke andere
oplossing verschilt van deze met een veelvoud van $\mathrm e^{-t/\tau}$, dat uitdooft. De
vergelijking is lineair, dus oplossingen tellen op.
\end{proof}

\begin{omfigure}
\centering
% figdata/bachelor-2/orbital.py part lag
\begin{tikzpicture}
\begin{axis}[width=0.96\linewidth, height=5.4cm, xmin=0, xmax=400, ymin=-3.4, ymax=4.2,
  xtick={0,50,...,400}, xlabel={ouderdom (duizenden jaren vóór 1950)},
  ylabel={gestandaardiseerd}, tick label style={font=\scriptsize}, label style={font=\small},
  legend style={font=\scriptsize, at={(0.5,1.02)}, anchor=south, draw=none, legend columns=3,
  /tikz/every even column/.append style={column sep=6pt}}, legend cell align=left]
\addplot[black!35, thin] table[x=age_ka, y=forcing] {figdata/out/bachelor-2-orbital-lag.dat};
\addlegendentry{min instraling \ang{65}~NB}
\addplot[ocRed, thick] table[x=age_ka, y=model] {figdata/out/bachelor-2-orbital-lag.dat};
\addlegendentry{vertragingsmodel, $\tau = \qty{10}{kyr}$}
\addplot[black, thick] table[x=age_ka, y=lr04] {figdata/out/bachelor-2-orbital-lag.dat};
\addlegendentry{LR04 (ijs omhoog)}
\end{axis}
\end{tikzpicture}
\par\smallskip{\small Het vertragingsmodel aangedreven door de instraling van half juli op \ang{65}~NB (minder
zon, meer ijs), tegenover de benthische stapelreeks, alle drie gestandaardiseerd (gemiddelde 0, standaardafwijking
1); positieve waarden betekenen meer ijs. Het model volgt de precessieschommelingen van de stapelreeks maar heeft
niets van haar grote zaagtandcycli. De responstijd van \qty{10}{kyr} is een modelkeuze.}
\end{omfigure}
% ledger: data:b4-orbital, data:lr04

Over de laatste \num{800000} jaar correleert het vertraagde model met de stapelreeks met $r = 0.42$, de
instraling zelf slechts met $0.23$: de vertraging helpt, maar een lineaire ijskap kan geen
cycli van \num{100000} jaar maken uit een forcering die er geen heeft. De terminaties moeten
niet-lineair gedrag inhouden dat een bescheiden duw van de instraling versterkt. De kandidaten zijn terugkoppelingen die op het hele klimaat werken:
de \omterm{def:b2:precambrian-earth:ice-albedo}{ijs--albedoterugkoppeling} van \cref{def:b2:precambrian-earth:ice-albedo}, want een planeet met meer
sneeuw en ijs weerkaatst meer zonlicht; en koolstofdioxide, dat in de Antarctische ijskernen stijgt op
tijdschalen van millennia telkens wanneer de omwentelende circulatie van de Atlantische Oceaan verzwakt, zoals
gebeurt wanneer smeltwater van de ijskappen in de Noord-Atlantische Oceaan stroomt. Voorbij een drempel van
koolstofdioxide kunnen interglaciale omstandigheden zich stabiliseren.
% ledger: data:b4-orbital, barker21:terminations, eo00:milankovitch

\begin{definition}[Midden-pleistocene overgang]\label{def:b2:orbital-forcing:mpt}
De \emph{midden-pleistocene overgang}\index{midden-pleistocene overgang} is de verandering, van ongeveer
\num{1250000} tot \num{700000} jaar geleden, van glaciale cycli van ongeveer \num{41000} jaar met kleine
amplitude naar quasi-periodieke cycli van ongeveer \num{100000} jaar met grote amplitude, met dikkere
ijskappen, zonder enige overeenkomstige verandering in de orbitale forcering.
\end{definition}
% ledger: clark06:mpt, nasa:milank-detail, kohler20:mpt

Vóór de overgang volgden de ijskappen de scheefstand, zoals het gestreepte spectrum van de figuur
hierboven toont; één telling vindt een \omterm{def:b2:orbital-forcing:glacials}{glaciale terminatie} in 67, 88 en \qty{52}{\percent} van de
scheefstandscycli van het vroege, middelste en late Kwartair. De orbitale theorie alleen kan de verandering niet verklaren; de kandidaten zijn de dynamica
van de ijskappen zelf, misschien na het verdwijnen van de zachte regoliet waarover ze gleden,
en een geleidelijke daling van koolstofdioxide tijdens de glaciale minima. Er is nog geen verklaring
die algemeen wordt aanvaard.
% ledger: kohler20:mpt, nasa:milank-detail, clark06:mpt

\section{Koolstofdioxide in de glaciale cycli}\label{sec:b2:orbital-forcing:co2}

De Antarctische ijskernen tonen dat koolstofdioxide de glaciale cycli volgt: tussen 188 en
\qty{194}{ppm} tijdens het Laatste Glaciale Maximum, tussen 266 en \qty{282}{ppm} in het laatste \omterm{def:b2:orbital-forcing:glacials}{interglaciaal};
de samengestelde reeks van \cref{ch:b2:climate-proxies} geeft \qty{189.7}{ppm} voor 19 tot
\qty{23}{ka} en \qty{279.4}{ppm} over de tweeduizend jaar vóór de industrie. Zijn stralingsforcering
groeit ongeveer met de logaritme van de concentratie, \qty{3.93}{W.m^{-2}} per verdubbeling.
% ledger: ipcc:fig2-34, data:b4-proxies-co2, ipcc:log-co2, ipcc:erf-2xco2

\begin{example}[Baan tegenover koolstofdioxide]\label{ex:b2:orbital-forcing:co2}
Van \qty{189.7}{ppm} naar \qty{279.4}{ppm} is de forcering $3.93\log_2(279.4/189.7) = 3.93 \times
0.559 = \qty{2.2}{W.m^{-2}}$, over de hele aarde en het hele jaar. De baan veranderde over het laatste
miljoen jaar het globaal en jaarlijks gemiddelde van de instraling met hoogstens $340.54 - 340.00 =
\qty{0.54}{W.m^{-2}}$, maar de instraling van half juli op \ang{65}~NB met ongeveer \qty{109}{W.m^{-2}}.
\end{example}
% ledger: data:b4-proxies-co2, ipcc:erf-2xco2, ipcc:log-co2, data:b4-orbital

De baan is een gangmaker, geen motor. Ze herverdeelt het zonlicht tussen seizoenen en breedten
en bepaalt het tijdstip; de ijskappen, door hun albedo en hun traagheid, en koolstofdioxide, door
zijn globale forcering, maken van een regionaal signaal een globale klimaatverandering.
% ledger: hays76:periods, barker21:terminations, eo00:milankovitch

\section*{Oefeningen}\label{sec:b2:orbital-forcing:exercises}

Gegevens voor de oefeningen, tenzij anders vermeld: $S_0 = \qty{1360}{W.m^{-2}}$ (de waarde van de
tabellen); een tropisch jaar van \num{365.2422} dagen; de baanelementen van de tabellen, waarvan vandaag
$e = 0.017236$, $\varpi = \ang{101.37}$, $\varepsilon = \ang{23.446}$.
% ledger: berger91:readme, berger91:rows

\begin{exercise}[$\star$]\label{exo:b2:orbital-forcing:1}
Neem $e = 0.0167$. (a) Bereken de verhouding van de apheliumafstand $a(1 + e)$ tot de
periheliumafstand $a(1 - e)$. (b) Druk hun verschil uit als percentage van $a$. (c) Met welke factor is
het zonlicht in het perihelium sterker dan in het aphelium?
\end{exercise}
% ledger: earth:ecc, nasa:milank-detail

\begin{exercise}[$\star$]\label{exo:b2:orbital-forcing:2}
De tabellen geven vandaag $e = 0.017236$, $\varpi = \ang{101.37}$; $e = 0.019971$, $\varpi =
\ang{293.86}$ \num{10000} jaar geleden; $e = 0.019249$, $\varpi = \ang{359.99}$ \num{6000} jaar geleden.
(a) Bereken de \omterm{def:b2:orbital-forcing:elements}{klimatologische precessie} voor elke ouderdom. (b) Geef voor elke ouderdom de lengte van de zon in het
perihelium. (c) Op welke ouderdom lag het perihelium het dichtst bij de junizonnewende, en op welke bij de
septemberequinox?
\end{exercise}
% ledger: berger91:rows

\begin{exercise}[$\star$]\label{exo:b2:orbital-forcing:3}
De rotatieas precedeert met een periode van ongeveer \num{25770} jaar ten opzichte van de sterren, terwijl
de baan draait met een periode van ongeveer \num{112000} jaar. Neem aan dat beide bewegingen het
perihelium ten opzichte van het lentepunt in dezelfde zin verplaatsen, zodat hun frequenties optellen. (a) Combineer
de twee frequenties om de periode te vinden waarmee het perihelium
door de seizoenen beweegt. (b) Vergelijk met de hoofdperioden van de \omterm{def:b2:orbital-forcing:elements}{klimatologische precessie}, \num{19.0},
\num{22.4} en \num{23.7} duizend jaar.
\end{exercise}
% ledger: nasa:milank-detail, la04:precession

\begin{exercise}[$\star$]\label{exo:b2:orbital-forcing:4}
Bereken met \cref{prop:b2:orbital-forcing:annual} en $e = 0.0167$ de jaargemiddelde instraling op
de noordpool voor scheefstanden van \ang{22.1} en \ang{24.5}. Met welk percentage groeit ze tussen
de twee?
\end{exercise}
% ledger: nasa:milank-detail, earth:ecc

\begin{exercise}[$\star\star$]\label{exo:b2:orbital-forcing:5}
(a) Bereken exact de afstandsfactor $(a/r)^2$ bij de junizonnewende vandaag en \num{116000} jaar
geleden ($e = 0.043988$, $\varpi = \ang{92.71}$), en vergelijk met de benadering $1 -
2e\sin\varpi$. (b) Bereken beide bij de decemberzonnewende. (c) Waarom is de benadering \num{116000}
jaar geleden slechter?
\end{exercise}
% ledger: berger91:rows

\begin{exercise}[$\star\star$]\label{exo:b2:orbital-forcing:6}
Werk \cref{ex:b2:orbital-forcing:today} stap voor stap opnieuw uit: de uurhoek van de zonsondergang op \ang{65}~NB
bij de junizonnewende, de daglengte, de instraling op de gemiddelde afstand en de
instraling van vandaag.
\end{exercise}
% ledger: berger91:rows

\begin{exercise}[$\star\star$]\label{exo:b2:orbital-forcing:7}
(a) Toon met $\mathrm dt = T(r/a)^2\,\mathrm d\lambda/(2\pi\sqrt{1 - e^2})$ en $(r/a)^2 \approx 1 +
2e\cos(\lambda - \varpi)$ aan dat het halfjaar van de maart- tot de septemberequinox
ongeveer $T/2 + (2T/\pi)\,e\sin\varpi$ duurt. (b) Bereken hoeveel dagen het vandaag langer is dan het andere halfjaar.
(c) Vergelijk met de 4.5 dagen waarmee noordelijke zomers langer zijn dan winters en de drie dagen waarmee
lentes langer zijn dan herfsten.
\end{exercise}
% ledger: nasa:milank-detail, berger91:rows

\begin{exercise}[$\star\star$]\label{exo:b2:orbital-forcing:8}
Over het laatste miljoen jaar varieerde de excentriciteit van \num{0.0046} tot \num{0.0565}. (a) Bereken
de globaal en jaarlijks gemiddelde instraling $S_0/(4\sqrt{1 - e^2})$ bij beide uitersten. (b) Hoeveel, in
\unit{W.m^{-2}} en in procent, veranderde de baan haar? (c) Vergelijk met het bereik van de
instraling van half juli op \ang{65}~NB, 388 tot \qty{497}{W.m^{-2}}.
\end{exercise}
% ledger: data:b4-orbital

\begin{exercise}[$\star\star$]\label{exo:b2:orbital-forcing:9}
Bereken met het vertragingsmodel en $\tau = \qty{10}{kyr}$ de relatieve amplitude en de vertraging van het
ijsvolume voor een sinusvormige forcering met perioden van 23, 41 en \num{100} duizend jaar. Welke band
begunstigt de ijskap, en waarom?
\end{exercise}

\begin{exercise}[$\star\star\star$]\label{exo:b2:orbital-forcing:10}
De instraling van half juli op \ang{65}~NB heeft bijna geen vermogen bij \num{100000} jaar, en toch heeft het ijsvolume
van de laatste \num{700000} jaar daar het grootste deel van zijn vermogen. (a) Toon aan dat als de forcering $A_1\sin
\omega_1 t + A_2\sin\omega_2 t$ is, het vertragingsmodel alleen bij $\omega_1$ en $\omega_2$ reageert. (b) Toon aan
dat een respons met de term $V^2$, bijvoorbeeld $\mathrm dV/\mathrm dt = (V_{\mathrm{eq}} -
V)/\tau + kV^2$, de frequenties $\omega_1 - \omega_2$ en $\omega_1 + \omega_2$ kan scheppen, door
de lineaire oplossing in $V^2$ in te vullen. (c) Met welke periode zweven twee precessietermen met perioden van \num{23.7} en
\num{19.0} duizend jaar? Wat volgt daaruit voor de cycli van \num{100000}
jaar?
\end{exercise}
% ledger: la04:precession, barker21:terminations

\begin{exercise}[$\star\star\star$]\label{exo:b2:orbital-forcing:11}
Op de noordpool is de energie die per vierkante meter tijdens de pooldag wordt ontvangen $E = \int S_0(a/r)^2
\sin\delta\,\mathrm dt$ over het halfjaar waarin $\delta > 0$. (a) Toon met het bewijs van
\cref{prop:b2:orbital-forcing:annual} aan dat $E = S_0T\sin\varepsilon/(\pi\sqrt{1 - e^2})$. (b) Bereken
$E$ in \unit{GJ.m^{-2}} voor $\varepsilon = \ang{22.1}$ en \ang{24.5}, met $e = 0.0167$ en $T$ =
365.25 dagen. (c) De precessie komt niet voor in $E$. Waarom kan ze toch het
afsmelten van een ijskap veranderen?
\end{exercise}
% ledger: nasa:milank-detail

\begin{exercise}[$\star\star\star$]\label{exo:b2:orbital-forcing:12}
(a) Bereken de stralingsforcering van de stijging van koolstofdioxide van \qty{189.7}{ppm} naar
\qty{279.4}{ppm}, met \qty{3.93}{W.m^{-2}} per verdubbeling en een logaritmische wet. (b) De bereiken
die voor het Laatste Glaciale Maximum en het laatste \omterm{def:b2:orbital-forcing:glacials}{interglaciaal} worden gegeven, zijn 188 tot \qty{194}{ppm} en 266 tot
\qty{282}{ppm}: bereken de kleinste en de grootste forcering die ze toelaten. (c) Vergelijk met de verandering
van de globale jaarlijkse instraling van \cref{exo:b2:orbital-forcing:8} en met de schommeling van de
zomerinstraling op \ang{65}~NB. Wat is de trekker en wat de versterker?
\end{exercise}
% ledger: data:b4-proxies-co2, ipcc:fig2-34, ipcc:erf-2xco2, ipcc:log-co2, data:b4-orbital

\begin{problem}[{Weekendprobleem --- Zonlicht op \ang{65}~NB}]\label{pb:b2:orbital-forcing:1}
Het laatste \omterm{def:b2:orbital-forcing:glacials}{interglaciaal} eindigde ongeveer \num{116000} jaar geleden. Hoeveel zonlicht ontving de zomer op de
breedte van de grote noordelijke ijskappen toen, vergeleken met nu? Gebruik de gegevens van de
oefeningen. De tabellen geven, \num{116000} jaar geleden, $e = 0.0440$, $\varpi = \ang{92.7}$,
$\varepsilon = \ang{22.52}$; vandaag, op dezelfde manier afgerond, $e = 0.0172$, $\varpi = \ang{101.4}$,
$\varepsilon = \ang{23.45}$.

\medskip\noindent\textbf{Deel I --- De baan toen en nu.}
\begin{enumerate}
  \item Bereken de \omterm{def:b2:orbital-forcing:elements}{klimatologische precessie} toen en nu.
  \item Bereken de perihelium- en apheliumafstanden, in eenheden van $a$, toen en nu.
  \item Bij welke lengte van de zon stond de aarde toen in het perihelium, en nu? Bij welke zonnewende?
  \item Met welke factor was het zonlicht in het perihelium sterker dan in het aphelium, toen en nu?
  \item Was de scheefstand groter of kleiner dan nu? Wat betekent dat voor de zomers op
        hoge breedten?
  \item Bereken de globaal en jaarlijks gemiddelde instraling toen en nu. Is het verschil significant?
\end{enumerate}

\medskip\noindent\textbf{Deel II --- De afstand bij de junizonnewende.}
\begin{enumerate}[resume]
  \item Toon met \cref{prop:b2:orbital-forcing:distance} aan dat bij de junizonnewende $r/a = (1 -
        e^2)/(1 - e\sin\varpi)$.
  \item Bereken $r/a$ bij de junizonnewende toen en nu.
  \item Bereken $(a/r)^2$ toen en nu.
  \item Vergelijk met de benadering $1 - 2e\sin\varpi$.
  \item Met welk percentage was het zonlicht bij de zonnewende zwakker dan op de gemiddelde afstand, toen en nu?
\end{enumerate}

\medskip\noindent\textbf{Deel III --- Zonlicht op \ang{65}~NB.}
\begin{enumerate}[resume]
  \item Geef de declinatie van de zon bij de junizonnewende toen en nu.
  \item Bereken de uurhoek van de zonsondergang op \ang{65}~NB op die dag, toen en nu.
  \item Leid daaruit de daglengte af.
  \item Bereken de daggemiddelde instraling op de gemiddelde afstand.
  \item Bereken de daggemiddelde instraling, toen en nu.
  \item De tabellen geven voor half juli 397.60 en \qty{426.76}{W.m^{-2}}. Waarom zijn de waarden
        bij de zonnewende hoger?
  \item Wat telde zwaarder in het verschil tussen toen en nu, de scheefstand of de afstand?
\end{enumerate}

\medskip\noindent\textbf{Deel IV --- Het antwoord van de ijskap.}
\begin{enumerate}[resume]
  \item Neem het vertragingsmodel met $\tau = \qty{10}{kyr}$: welk deel van een plotselinge verandering van
        $V_{\mathrm{eq}}$ heeft de ijskap bereikt na 5, 10 en \qty{20}{kyr}?
  \item De juni-instraling op \ang{65}~NB daalde van een maximum ongeveer \num{128000} jaar geleden tot het
        minimum van \num{116000} jaar geleden, een halve precessiecyclus. Welk deel van de forcering volgt de ijskap
        voor een sinusoïde met periode \qty{23}{kyr}, en met hoeveel vertraging?
  \item Waarom kan een ijskap binnen één precessiecyclus niet de grootte bereiken die een lage
        zomerinstraling in stand zou houden?
  \item Het ijs had zo'n \num{90000} jaar nodig om op te bouwen tot het Laatste Glaciale Maximum en ongeveer
        \num{10000} jaar om af te smelten. Waarom kan het vertragingsmodel die asymmetrie niet voortbrengen?
  \item Koolstofdioxide stond tussen 266 en \qty{282}{ppm} in het laatste \omterm{def:b2:orbital-forcing:glacials}{interglaciaal} en tussen
        188 en \qty{194}{ppm} tijdens het Laatste Glaciale Maximum. Hoe helpt het een regionale verandering van
        zonlicht een globale klimaatverandering te worden?
  \item Samengevat: wat was de daggemiddelde instraling op \ang{65}~NB bij de junizonnewende toen het
        laatste \omterm{def:b2:orbital-forcing:glacials}{interglaciaal} eindigde, \num{116000} jaar geleden?
\end{enumerate}
\end{problem}
% ledger: berger91:rows, berger91:readme, data:b4-orbital, ipcc:ref-periods, ipcc:fig2-34, barker21:terminations
